% Gesamttest «Bruchterme, Gleichungen und Formeln» — Quelle des PDFs (python3 scripts/build-lp-pdf.py gleichungen-formeln).
% Musterlösung und Raster: bewertungspaket.tex. Alle Werte mit python3 nachgerechnet (scripts/lp/gleichungen-formeln/zahlen.py).
% Selbstcheck nach ~/TODO-BM-Vorkurs-LP03-15.md: Bruchterm, Gleichung, Formel (Treppenmodell als Prüfungsanker).
% 25 Punkte, rund 30 Minuten. Nach der Prüfung vom 10.10.2026: G2 mit Division (K2 ganz geprüft), G3 2 P; Kapitelaufgabe 4f und
% die Übung «Treppe» ersetzt, damit G7 nichts wiederholt.
\documentclass[11pt]{article}
\usepackage{../lp-druck}
\renewcommand{\lpname}{Bruchterme, Gleichungen und Formeln}
\renewcommand{\lpteilgebiet}{Gleichungen ·}
\renewcommand{\lpfuss}{Gesamttest · 25 Punkte}
\begin{document}
\lpkopf{Gesamttest}{%
  \vspace{4pt}\noindent\begin{tabularx}{\linewidth}{@{}X X@{}}
  Code (kein Name): \hrulefill & Punkte: \hrulefill\ / 25
  \end{tabularx}}

\begin{lpinfo}{Anleitung}
\textbf{Zeit:} rund 30 Minuten. \textbf{Hilfsmittel:} keine, kein Taschenrechner.
Schreib zu jeder Aufgabe den \textbf{Rechenweg} auf; bewertet wird der Weg. Gib Bruchterme \textbf{vollständig gekürzt} auf einem Bruchstrich an,
Lösungen als \textbf{Lösungsmenge}.
Danach: selbst mit dem \emph{Bewertungspaket} bewerten oder die Lösung samt diesem Blatt scannen bzw.\ fotografieren
(ohne Namen und Standort) und mit dem Bewertungspaket einer KI geben.
\end{lpinfo}

\lpteil{Teil A · Bruchterme}{Kapitel 1 und 2 · 9 Punkte}

\begin{lpaufgabe}{G1}{Kürzen}{3 P}
Kürze vollständig: $\;\dfrac{x^2 - 10x + 25}{3x - 15}$. Welche Zahl darf $x$ nicht sein?
\lpplatz{5}
\end{lpaufgabe}

\begin{lpaufgabe}{G2}{Rechnen}{4 P}
\begin{enumerate}[label=(\alph*),leftmargin=*,itemsep=4pt]
\item Dividiere und kürze vollständig: $\;\dfrac{x^2 - 4}{3x} : \dfrac{x + 2}{6x^2}$
\item Addiere und schreib das Ergebnis auf einen Bruchstrich: $\;\dfrac{3}{x} + \dfrac{2}{x - 1}$
\end{enumerate}
\lpplatz{7}
\end{lpaufgabe}

\begin{lpaufgabe}{G3}{Fehler finden}{2 P}
Jemand kürzt: $\;\dfrac{x^2 + 9}{x + 3} = x + 3$. Zeige mit $x = 1$, dass das nicht stimmt.
Lässt sich der Bruch überhaupt kürzen? Begründe.
\lpplatz{5}
\end{lpaufgabe}

\lpteil{Teil B · Gleichungen}{Kapitel 3 · 8 Punkte}

\begin{lpaufgabe}{G4}{Klammern}{3 P}
Löse und mach die Probe: $\;7x - 3(x - 2) = 2x + 14$
\lpplatz{5}
\end{lpaufgabe}

\begin{lpaufgabe}{G5}{Brüche}{3 P}
Löse: $\;\dfrac{x + 3}{4} - \dfrac{x - 1}{6} = 2$
\lpplatz{5}
\end{lpaufgabe}

\begin{lpaufgabe}{G6}{Lösungsmenge}{2 P}
Gib die Lösungsmenge an und begründe: $\;4(x - 1) + 6 = 4x + 2$
\lpplatz{3}
\end{lpaufgabe}

\lpteil{Teil C · Formeln}{Kapitel 4 · 8 Punkte}

\begin{lpaufgabe}{G7}{Treppe}{3 P}
Für bequeme Treppen gilt die Schrittmassregel $\;2h + a = 63\;$ (in cm; $h$ Stufenhöhe, $a$ Auftrittsbreite).
\begin{enumerate}[label=(\alph*),leftmargin=*,itemsep=3pt]
\item Stelle die Formel nach $h$ um.
\item Wie hoch ist eine Stufe mit der Auftrittsbreite $a = 25$ cm?
\end{enumerate}
\lpplatz{5}
\end{lpaufgabe}

\begin{lpaufgabe}{G8}{Zins}{3 P}
Für den Zins ohne Zinseszins gilt $\;Z = \dfrac{K \cdot p \cdot t}{100}\;$ ($K$ Kapital in Fr., $p$ Zinssatz in Prozent, $t$ Jahre, $Z$ Zins in Fr.).
Stelle nach $p$ um. Welcher Zinssatz bringt auf $2000$ Fr. in $3$ Jahren $150$ Fr. Zins?
\lpplatz{5}
\end{lpaufgabe}

\begin{lpaufgabe}{G9}{Fehler finden}{2 P}
Pyramide: $\;V = \dfrac{G \cdot h}{3}$. Jemand stellt nach $G$ um und schreibt $\;G = \dfrac{V \cdot h}{3}$.
Prüfe mit $G = 12$ und $h = 5$, und stelle richtig um.
\lpplatz{4}
\end{lpaufgabe}
\end{document}
